\begin{center}\footnotesize
\begin{longtable}{rlp{7.4cm}p{4.6cm}}
\hline
Ch. & Paper & Title & Source directory\\
\hline
\endhead
1 & --- & KMS state uniqueness and phase transitions in localized Bost--Connes systems:  from Chebotarev density to the Hardy--Littlewood conjecture & \nolinkurl{paper/}\\
2 & I & Factor type and the spectrum of transitions & \nolinkurl{papers/I-type-spectrum/}\\
3 & II & Ideal structure, boundary quotient, and index & \nolinkurl{papers/II-structure/}\\
4 & III & Invariants, entropy, and the measure class & \nolinkurl{papers/III-invariants/}\\
5 & XXXIX & The measurable first cohomology of the tail relation:  turbulence and the arithmetic embedding & \nolinkurl{papers/XXXIX-cohomology/}\\
6 & IV & Noncommutative metric geometry & \nolinkurl{papers/IV-metric/}\\
7 & XXIX & The Wasserstein separation of the broken phases:  a finite-dimensional reduction and the quadratic chaos bound & \nolinkurl{papers/XXIX-wasserstein/}\\
8 & XXXIII & The adaptive coupling upper bound for the Wasserstein separation:  microscopic flips against the macroscopic symmetry cost & \nolinkurl{papers/XXXIII-coupling/}\\
9 & XXXVII & The microscopic geometry of Wasserstein separation:  cone geodesics, the exact single-step barrier, and the failure of fractional transport & \nolinkurl{papers/XXXVII-cone-geometry/}\\
10 & XXX & K-theory of the affine boundary quotient:  the return of the norm and the trace & \nolinkurl{papers/XXX-affine-ktheory/}\\
11 & XXXI & The ray-class repair for affine systems:  the class group as the index of definable scalings & \nolinkurl{papers/XXXI-rayclass/}\\
12 & XXXII & Galois descent for affine Bost--Connes boundaries:  capitulation, the obstruction cocycle, and the emergence of the class group & \nolinkurl{papers/XXXII-descent/}\\
13 & XXXVIII & Integral equivariant K-theory of the inert local model:  the ray-class entanglement and exact torsion & \nolinkurl{papers/XXXVIII-inert-ktheory/}\\
14 & V & The free variant and the free critical exponent & \nolinkurl{papers/V-free/}\\
15 & VI & Towards non-abelian localized Bost--Connes systems:  the matrix Kakutani dichotomy and a Tannakian obstruction & \nolinkurl{papers/VI-nonabelian/}\\
16 & VII & The Galois-twisted arithmetic groupoid:  its KMS classification, and why the Tannakian obstruction  does not compute it & \nolinkurl{papers/VII-galois-groupoid/}\\
17 & XXVII & Non-abelian Bost--Connes groupoids:  two no-go theorems, and a construction that realizes $\Prob(G/N_\beta)$  exactly when the Frobenius lifts mix & \nolinkurl{papers/XXVII-nonabelian-groupoids/}\\
18 & XXXIV & The abelian scaling bottleneck, and the angular Hurwitz system:  where noncommutative arithmetic thermodynamics can and cannot live & \nolinkurl{papers/XXXIV-bottleneck/}\\
19 & XXXVI & The angular transport: tree-level filtrations, predictable matrix martingales,  and the two regimes of the angular Hurwitz system & \nolinkurl{papers/XXXVI-angular-transport/}\\
20 & IX & Localized Bost--Connes factors are amenable:  the failure of $W^*$-superrigidity,  and what survives at the $C^*$-level & \nolinkurl{papers/IX-rigidity/}\\
21 & XIII & Equivariant rigidity of localized Bost--Connes  $C^*$-dynamical systems: the fixed-point algebra  is the Cartan subalgebra & \nolinkurl{papers/XIII-equivariant/}\\
22 & XIV & The equal-norm case: norm-separation is not a hypothesis  but a characterization, and the gauge-equivariant repair & \nolinkurl{papers/XIV-equalnorm/}\\
23 & XV & What $\sigma$-equivariance gives for free,  and why the range projections cannot give the rest & \nolinkurl{papers/XV-spectral/}\\
24 & XVI & Quantum optimal transport on Bost--Connes systems:  the twist is harmful, the generator has a sign error,  and the transition is invisible to local inequalities & \nolinkurl{papers/XVI-transport/}\\
25 & X & A non-commutative Ore LCM semigroup  with arithmetic thermodynamics:  the Hurwitz quaternions & \nolinkurl{papers/X-semigroup-search/}\\
26 & XI & Condition \textup{(C4)} and the nested/partition dichotomy:  why Ramanujan bounds cannot decide  non-amenability of the Hurwitz groupoid & \nolinkurl{papers/XI-C4/}\\
27 & XII & The universal multiplicity--norm mismatch,  and the reduction of \textup{(C4)} to the norm-one subgroup & \nolinkurl{papers/XII-mismatch/}\\
28 & XXI & Localized $GL_2$ Bost--Connes systems:  the equilibrium range and the automorphic obstruction range  are disjoint, by exactly one & \nolinkurl{papers/XXI-GL2/}\\
29 & XXVIII & The Ramanujan bound is the tempered bound:  the Hecke operator on the Bruhat--Tits tree, and  a no-go theorem for condition (C4) & \nolinkurl{papers/XXVIII-spectral-gap/}\\
30 & VIII & Localized Bost--Connes systems in arithmetic topology:  local class field theory, linking Kakutani dichotomies,  and topological phase transitions & \nolinkurl{papers/VIII-arithmetic-topology/}\\
31 & XVII & Localized Bost--Connes systems over global function fields:  commensurability forbids type $\mathrm{III}_1$,  and the factor recovers the temperature & \nolinkurl{papers/XVII-function-fields/}\\
32 & XVIII & Circle actions on function field Bost--Connes systems:  the compact group advantage is illusory,  and the Weil zeta function is thermodynamic & \nolinkurl{papers/XVIII-equivariant-kirchberg/}\\
33 & XIX & Isogeny, zeta-equivalence, and the Picard torsor  in function field Bost--Connes systems & \nolinkurl{papers/XIX-isogeny/}\\
34 & XX & The limits of abelian reconstruction:  twin curves exist, genus one suffices,  and the obstruction is a torsor coordinate & \nolinkurl{papers/XX-twin-curves/}\\
35 & XXII & Iwasawa theory and localized Bost--Connes systems:  there are no $p$-adic KMS states,  and where Leopoldt actually enters & \nolinkurl{papers/XXII-iwasawa/}\\
36 & XXIII & Leopoldt and the $p$-local Bost--Connes groupoid:  the rank of the isotropy, not its triviality,  and why $\log_p(\Ns\pp)$ vanishes & \nolinkurl{papers/XXIII-leopoldt/}\\
37 & XXIV & The statistical mechanics of sieve-theoretic primes:  the critical spectrum, the log-power dichotomy,  and two families that are never unique & \nolinkurl{papers/XXIV-sieve/}\\
38 & XXV & Sato--Tate families in localized Bost--Connes systems:  open phase, no congruence obstruction,  and the symmetric power expansion & \nolinkurl{papers/XXV-sato-tate/}\\
39 & XXXV & The free twin-prime Bost--Connes system is unconditionally supercritical:  resolving the phase dichotomy via exact partial sums & \nolinkurl{papers/XXXV-twin-regime/}\\
\hline
\end{longtable}
\end{center}
